\section{Conclusion}
As machine learning is becoming a major tool for researchers and policy makers across different fields such as healthcare and economics, causal inference becomes a crucial issue for the practice of machine learning. In this paper we focus on counterfactual inference, which is a widely applicable special case of causal inference. We cast counterfactual inference as a type of domain adaptation problem, and derive a novel way of learning representations suited for this problem.

Our models rely on a novel type of regularization criteria: learning \emph{balanced representations}, representations which have similar distributions among the treated and untreated populations. We show that trading off a balancing criterion with standard data fitting and regularization terms is both practically and theoretically prudent.

Open questions which remain are how to generalize this method for cases where more than one treatment is in question, deriving better optimization algorithms and using richer discrepancy measures.
